\chapter{Metrados y presupuesto}
\thispagestyle{memoria}\pagestyle{memoria}

\section{Base del metrado}
Los metrados se generan desde las longitudes de los 35 circuitos y cinco
alimentadores del modelo electrico, mas el alimentador general y las cantidades
de alumbrado, tableros, respaldo, puesta a tierra y dispositivos declaradas en
los datos canonicos. Se aplica 10\% de holgura a rutas y conductores. Las
longitudes deben contrastarse en campo antes de comprar.

\section{Alcance economico}
El presupuesto tiene fecha base 2 de agosto de 2026 y estado
\textbf{\EstadoPresupuesto}. Es una estimacion academica reproducible, no una
cotizacion. Los precios visibles en tiendas peruanas solo sirven como referencias
comparables; los equipos industriales y trabajos especializados requieren
proformas formales. Se excluyen surtidores, STP, tanques, tuberias de combustible,
licencias, derechos de conexion y obra civil general.

El costo directo calculado es \textbf{S/ \CostoDirecto} y el total referencial,
incluidos gastos generales, utilidad e IGV, es \textbf{S/ \PresupuestoTotal}.

\section{Cuadro de metrados y presupuesto}
\footnotesize
\input{../../../build/unidad-2-industrial/expediente/generated/metrados-presupuesto}
\normalsize

\section{Regla de actualizacion}
Antes de una compra real se deben reemplazar todas las filas
\texttt{estimado\_anteproyecto} por cotizaciones vigentes, confirmar marcas y
certificaciones, recalcular cantidades con recorrido de obra y registrar flete a
Caracoto. Una diferencia de precio no autoriza reducir seccion, poder de corte,
grado de proteccion ni certificacion para area clasificada.
